% ============================================================
% KBS — Knowledge-Based Systems Cover Letter
% Elsevier elsarticle
% ============================================================

\documentclass[12pt]{article}
\usepackage[T1]{fontenc}
\usepackage[utf8]{inputenc}
\usepackage[margin=1in]{geometry}
\usepackage{xurl}

\begin{document}

\begin{center}
{\Large \textbf{Cover Letter — Knowledge-Based Systems}}
\end{center}

\vspace{12pt}

\textbf{Date:} July 28, 2026

\vspace{6pt}

\textbf{To:} Prof. Jie Lu, Editor-in-Chief \\
\textit{Knowledge-Based Systems} \\
Elsevier

\vspace{12pt}

Dear Prof. Lu,

I am pleased to submit our manuscript entitled \textbf{``A Knowledge-Based System for Educational Assessment: Knowledge Engineering, Multi-Method Explainable AI, and Fairness-Aware Decision Support on PISA 2022 Data''} for consideration in \textit{Knowledge-Based Systems}.

\section*{Fit with Knowledge-Based Systems}

This manuscript presents a knowledge-based system that encodes domain expertise from educational research and digital inequality scholarship into a reproducible machine learning workflow for large-scale educational assessment. The work aligns with KBS's focus on knowledge-based and other artificial intelligence techniques-based systems that support human prediction and decision-making, in the following ways:

\begin{enumerate}
\item \textbf{Knowledge engineering:} The system's knowledge engineering layer organizes 33 PISA 2022 predictors (613,744 students, 80 countries/economies) using ecological systems theory and a multi-level ICT taxonomy. This is not merely feature engineering---it is the systematic encoding of domain knowledge from educational psychology and digital divide research into a structured analytical framework.

\item \textbf{Structured knowledge validation:} Six domain-knowledge-derived propositions (P1--P6) are mapped to empirical tests, operationalizing theoretical claims as testable predictions against fitted-model behavior (Table 11). This constitutes a knowledge validation protocol that makes domain assumptions explicit and empirically checkable.

\item \textbf{Multi-method explanation with cross-validation:} SHAP, permutation importance, and Accumulated Local Effects converge (Spearman's $\rho = 0.76$--$0.83$), while LIME diverges substantially ($\rho < 0.03$), providing a structured comparison of explanation methods that informs method selection for correlated tabular data.

\item \textbf{Fairness-aware decision support:} Formal fairness metrics (Equalized Odds, Demographic Parity, ABROCA) with intersectional subgroup analysis establish quantitative baselines for equitable deployment.

\item \item \textbf{Reproducibility:} The complete pipeline operates on a single workstation (Apple M5, $\sim$12 hours) with fixed random seeds and deterministic explanation samples, with code publicly available.

\item \textbf{Engagement with KBS literature:} The revised manuscript directly engages with recent KBS publications, including Cheng et al.\ (2026, uncertainty explanation by SHAP), Wang and Luo (2024, latent discrimination detection), and Yang et al.\ (2025, integrating LLMs with knowledge-based methods), in addition to earlier KBS papers by Zhang et al.\ (2025, fairness-aware feature selection) and Pei (2025, fair federated learning). Six of 59 references (10.2\%) are now from Knowledge-Based Systems.
\end{enumerate}

\section*{Key Results}

\begin{itemize}
\item Tuned XGBoost: weighted AUC = 0.903, RMSE = 54.10 ($R^2 = 0.681$), 23\% improvement over linear baselines
\item Supplementary $k$-fold CV confirms negligible partition sensitivity (AUC = 0.890 $\pm$ 0.001)
\item SES disparities identified as largest fairness concern (ABROCA = 0.023)
\item All six knowledge-derived propositions received empirical support
\end{itemize}

\section*{Relevance Statement}

This manuscript was previously submitted to \textit{Expert Systems with Applications}, where editorial review indicated that the expert-system framing was not sufficiently central. We recognize that the work's core contribution is not an expert system in the deployed-production sense, but rather a knowledge-based system in the knowledge-engineering sense: domain knowledge (ecological systems theory, ICT taxonomy) is systematically encoded to constrain and guide the ML workflow from feature selection through interpretation, explanation, and fairness evaluation---which aligns directly with KBS's core identity. We have substantially revised the manuscript to reflect this framing throughout, including renaming the five system layers, adding a knowledge encoding protocol, introducing a dedicated knowledge validation sub-section, and reframing the multi-method XAI comparison as a knowledge extraction framework.

\section*{Declarations}

\textbf{Data and Code Availability.} PISA 2022 data are publicly available from the OECD. Analysis code and non-restricted derived outputs are available at \url{https://github.com/Jackxiaozhiren/pisa2022-xai-math}.

\textbf{Declaration of Competing Interest.} The author declares no competing interests.

\textbf{Declaration of Generative AI and AI-Assisted Technologies.} AI-assisted tools were used for programming support, reproducible analysis organization, literature search, and language editing. They were NOT used for study design, data interpretation, theoretical framework selection, drawing substantive conclusions, or final manuscript approval. This statement is consistent with Elsevier's AI authorship policy and COPE guidelines.

\textbf{Originality.} This manuscript is original, has not been published previously, and is not under consideration elsewhere. It is a single-author work; the author has read and approved the final version.

\textbf{Suggested Reviewers.}
\begin{enumerate}
\item Prof. Okan Bulut, University of Alberta --- large-scale educational assessment, machine learning
\item Prof. Crist\'{o}bal Romero, University of C\'{o}rdoba --- educational data mining
\item Prof. Ren\'{e} F. Kizilcec, Cornell University --- algorithmic fairness in education
\item Dr. Julia Herbinger, LMU Munich --- GADGET framework, feature effects
\item Dr. Hassan Khosravi, University of Queensland --- explainable AI in education
\end{enumerate}

\vspace{12pt}
Thank you for your time and consideration.

\vspace{18pt}

Sincerely,

\vspace{12pt}

\textbf{Zhiren Xiao} \\
Guangdong University of Finance \\
Email: 241734106@m.gduf.edu.cn \\
ORCID: 0009-0008-2164-4557

\end{document}
